\pagebreak

\chapter{Code structure}

    A program usually consists of modules interconnected by queues. A module
    code body contains one or more code elements each of which is initiated
    repeatedly, i.e. is contained in an infinite loop.

    A program file contains optional global scope and file scope structure and
    variable creation calls and executable code intermixed with module,
    procedure and function declarations. It is preferable to have
    top level code first followed by the module, procedure and function
    declarations.
    
    Execution begins with the initial file module which in turn calls other
    modules, procedures and functions.

    Execution of statements in {\bf 3PL} is controlled by two mechanisms -
    \blist
    \item    execution sequence
    \item    queue availability
    \blend
    
    Statements in a module (including the main file module) which are not nested
    within control structures, such as conditional statements or loops, will be
    refered to as {\bf top level} statements.

    All top level statements are executed in an implied infinite loop.
    Statements within control structures are executed under the control of those
    structures. Statements may additionally wait upon the availability of queues.

    Top level simple assignment statements occurring immediately within
    a module will execute whenever referenced queues are available.
    If no queues are referenced, the statements will execute on
    every clock cycle.

    An assignment statement enclosed in a {\bf when} which itself
    is immediately inside a module will be controlled primarily
    by the {\bf when} which will execute every clock cycle. The
    {\bf when} body may wait on queues when the {\bf when} is true,
    otherwise it will execute at once.

    A great deal of code can be written without using
    explicit serial and parallel constructs, relying on queue
    availability alone to control execution.

    A simple program might be -
    \begin{lstlisting}[gobble=4, language=threepl]
       cmemory(m, 2, 8, 8);
       queue(a, "uint:8");
       queue(b, "uint:8");
       queue(addr, "uint:8");
       queue(d, "bits:8");

       [addr <<- a + b;]
       [d <<- cmemoryread(d <- m, 0, addr);]
    \end{lstlisting}
    The two target statements within the initial file module are themselves
    modules. The first is an inline module and the single assignment
    statement within will execute when its queues are available. The second is
    also an inline module and again will wait on the availability of its
    queues.

    If a module has more than one statement, e.g. -
    \begin{lstlisting}[gobble=4, language=threepl]
       [
           y <<- a + b;
           z <<- c - d;
       ]
    \end{lstlisting}
    each statement in the module would execute when its
    queues were available. We could have simply made these
    independent modules -
    \begin{lstlisting}[gobble=4, language=threepl]
       [y <<- a + b;]
       [z <<- c - d;]
    \end{lstlisting}
    The first form would be used had we needed to share
    local static variables.

    There are three structures giving parallelism within a module
    which we must understand.

    The first has been used above, simply enclosing statements
    in inline module brackets
    \begin{lstlisting}[gobble=4, language=threepl]
       
          y <<- a + b;
          z <<- c - d;
       
    \end{lstlisting}
    Each statement in the module will execute independently and repeatedly
    whenever it can. If a statement contains only static variables it will
    execute on every clock cycle (such statements would normally be contained
    within a {\bf when} or other such conditional statement to control their
    execution). If a statement contains queue variables it will execute when
    these are available (both input on the RHS and output on the LHS).

    If we enclose the statements in a {\bf par} block
    \begin{lstlisting}[gobble=4, language=threepl]
       [
           par {
               y <<- a + b;
               z <<- c - d;
           }
       ]
    \end{lstlisting}
    each statement within the block will execute exactly once as
    soon as it can. The parallel block will then have completed
    execution and since it is immediately enclosed within a module
    the parallel block will execute repeatedly. Note that the
    statements within the parallel block are paired in execution but
    will not necessarily execute at the same time.

    If we wish to ensure that parallel statements execute simultaneously
    we must enclose them in a {\bf sync} block
    \begin{lstlisting}[gobble=4, language=threepl]
       [
           sync {
               y <<- a + b;
               z <<- c - d;
           }
       ]
    \end{lstlisting}
    which will execute both statements in the block when all queue variables
    in the block are available, i.e. they are guaranteed to execute
    simultaneously and in a single clock cycle. Note that if a {\bf sync}
    block contains a sequential block or an iterative statement the {\bf
    sync} simply controls entry into the block after which it behaves ilke a
    {\bf par} block, hence it would not control any further reads or
    writes of those queues.

    A situation to be wary of with queue dependencies occurs
    when queue fields are being assigned. For example
    \begin{lstlisting}[gobble=4, language=threepl]
       [
           v.a <<- y;
           v.b <<- z;
       ]
    \end{lstlisting}
    If {\bf y} and {\bf z} are not available simultaneously
    then the assignment will occur separately on each
    becoming available, the other field of {\bf v} being
    assigned zero. In that case they must be locked together by
    \begin{lstlisting}[gobble=4, language=threepl]
       [
           sync {
               v.a <<- y;
               v.b <<- z;
           }
       ]
    \end{lstlisting}        

    We can use parallel blocks to control static
    assignments, e.g.
    \begin{lstlisting}[gobble=4, language=threepl]
       [
           static(s, "uint:8");

           par {
               a <<- b + c;
               s++;
           }
       ]
    \end{lstlisting}
    where the static assignment {\bf s++} will only occur as
    often as the queue assignment, and the availability of
    input and output queues will control both. For actual
    simultaneous execution we would again write
    \begin{lstlisting}[gobble=4, language=threepl]
       [
           static(s, "uint:8");

           sync {
               a <<- b + c;
               s++;
           }
       ]
    \end{lstlisting}

    In the previous examples we might have been using {\bf s} to
    count queue data processed. If we wished to count total cycles
    to determine the proportion of active cycles we could write -
    \begin{lstlisting}[gobble=4, language=threepl]
       [
           static({s, t}, "uint:8");

           par {
               a <<- b + c;
               s++;
           }
           t++;
       ]
    \end{lstlisting}
    The static assignment to {\bf t} will occur every clock cycle
    since it is immediately within a module.

    For queues which are structures or arrays we can
    simply assign them, e.g.
    \begin{lstlisting}[gobble=4, language=threepl]
       struct {t1,
           d,  "uint:8",
           l,  "log"
       };
       queue({a, b}, t1);

       module
       m () {
           a <<- b;
       }
    \end{lstlisting}
    will assign the entire structure. We
    can mix different queue types as long as each component assignment
    type matches, e.g.
    \begin{lstlisting}[gobble=4, language=threepl]
       struct {t1,
           d,  "uint:8",
           l,  "log"
       };
       queue(a, t1);
       queue(pixel, "uint:8");
       queue(j, "log");

       module
       m () {
           sync {
               a.d <<- pixel;
               a.l <<- j;
           }
       }
    \end{lstlisting}
    but note that we have used the {\bf sync \{\dots\}} parallel
    block to ensure that the two fields are assigned simultaneously.
    If we do not ensure simultaneity then two separate data will
    we written, one with the {\bf l} field false, the other with
    the {\bf d} field zero. There is no conflict assigning to
    output queue {\bf a} as the assignments are to separate
    fields of the structure.

    If we want to examine two inputs but only advance one
    of them -
    \begin{lstlisting}[gobble=4, language=threepl]
       queue({o, i1, i2}, "uint:8");

       [
           sync {
               when (?i1 > ?i2)
                   o <<- i1;
               else
                   o <<- i2;
           }
       ]
    \end{lstlisting}
    The input queues {\bf i1} and {\bf i2} are examined
    without reading (the {\bf \verb'?'} operator) and then
    the larger is read and assigned to the output queue.
    The {\em when} will not wait for queues {\bf i1} and {\bf i2} to 
    become available in the test expression because the queues are examined,
    not read, but we have enclosed the {\em when} in a {\em sync} block
    and there are reads of queues {\bf i1} and {\bf i2} within that block so
    the {\em sync} will wait until queues {\bf i1}, {\bf i2} and {\bf o} are
    available before executing the {\em when}.

    We could shorten the code, using a {\bf? :} operator, to
    \begin{lstlisting}[gobble=4, language=threepl]
       queue({o, i1, i2}, "uint:8");

       [o <<-  (?i1 > ?i2) ? i1 : i2;]
    \end{lstlisting}
    Note that we now have a single assignment statement and
    that statement will not execute until all queues read or written
    (but not examined) within it are available.

    If we wanted to do the same but advance both inputs
    each time we could do this by
    \begin{lstlisting}[gobble=4, language=threepl]
       queue({o, i1, i2}, "uint:8");

       [o <<-  (i1 > i2) ? i1 : i2;]
    \end{lstlisting}
    since both inputs are read regardless of which is selected
    for assignment to the output queue.

    If we wished to take inputs when they become available
    and channel them to one output queue we might write -
    \begin{lstlisting}[gobble=4, language=threepl]
       // BAD CODE!
       queue({o, i1, i2}, "uint:8");

       [
           when (<i1)
               o <<- i1;
           when (<i2)
               o <<- i2;
       ]
    \end{lstlisting}
    but this is both pointless and incorrect. Firstly the
    availability tests are not necessary as the assignments will
    not execute unless the input queues are available anyway.
    Secondly if both input queues are available
    the assignments would conflict. We could check both
    explicitly by -
    \begin{lstlisting}[gobble=4, language=threepl]
       queue({o, i1, i2}, "uint:8");

       [
           when (<i1 && !<i2)
               o <<- i1;
           when (!<i1 && <i2)
               o <<- i2;
           when (<i1 && <i2)
               o <<- i1;
       ]
    \end{lstlisting}
    Here there are no conflicts as the three conditions are
    mutually exclusive (we do not need mutual exclusion
    like this if there is no assignment conflict).
    The above code will favour the first input when both are
    available. To give both inputs an even chance
    \begin{lstlisting}[gobble=4, language=threepl]
       queue({o, i1, i2}, "uint:8");
       static(first, "log");

       [
           when (<i1 && !<i2)
               o <<- i1;
           when (!<i1 && <i2)
               o <<- i2;
           when (<i1 && <i2) par {
               when  (first)
                   o <<- i1;
               else
                   o <<- i2;
               first <- !first;
           }
       ]
    \end{lstlisting}
    The inner {\em when} and the static assignment will be done in
    parallel in one cycle. It is preferable to not nest {\em when}s
    as they are not optimised and the nesting increases the length of logic
    paths. We could instead write -
    \begin{lstlisting}[gobble=4, language=threepl]
       queue({o, i1, i2}, "uint:8");
       static(first, "log");

       [
           when (<i1 && !<i2)
               o <<- i1;
           when (!<i1 && <i2)
               o <<- i2;
           when (<i1 && <i2 && first) par {
               o <<- i1;
               first <- false;
           }
           when (<i1 && <i2 && !first) par {
               o <<- i2;
               first <- true;
           }
       ]
    \end{lstlisting}

    The inner {\em when} could be contracted into a single assignment
    using the ternary operator {\bf? :} on the RHS.
    \begin{lstlisting}[gobble=4, language=threepl]
       queue({o, i1, i2}, "uint:8");
       static(first, "log");

       [
           if (<i1 && !<i2)
               o <<- i1;
           if (!<i1 && <i2)
               o <<- i2;
           if (<i1 && <i2) par {
               o <<- first ? i1 : i2;
               first <- !first;
           }
       ]
    \end{lstlisting}

    {\bf Value} variables serve two purposes - they can represent
    a combinatorial expression that is to be used in more than
    one place without duplicating the logic, and they can also
    be used simply to refer to other variables (the same thing
    really - an expression consisting of a single variable).

    The first usage is quite straightforward, for example
    \begin{lstlisting}[gobble=4, language=threepl]
       queue({q1, q2, q3, q4}, "uint:8");

       [
           q3 <<- q1 + q2;
           q4 <<- q1 + q2;
       ]
    \end{lstlisting}
    would duplicate the logic performing the addition, whereas
    \begin{lstlisting}[gobble=4, language=threepl]
       queue({a, b, c, d}, "uint:8");

       [
           value(v, "uint:8");

           v := q1 + q2;
           q3 <<- v;
           q4 <<- v;
       ]
    \end{lstlisting}
    would build only one instance of the addition. Note that where
    a value expression contains queues those queues will be read and
    acknowledged as though they appeared separately in the statements
    in which the value variable is used, even though any operators
    within the expression will only be instantiated once. This of course
    applies equally to queues appearing in array subscript expressions
    within a value expression.

    The second usage of a value variable is as a reference or pointer.
    For example we may wish to devise a module each invokation of which
    reports to a central module. We can do this as follows
    \begin{lstlisting}[gobble=4, language=threepl]
           queue({q1, q2, .... }, "uint:8");
           value(p_report, "[10]uint:8");
           int(reports);

           .
           .
           m(q1);
           .
           .
           m(q2);
           .
           .
           report();

       module
       m (queue(in, "uint:8")) {
           p_report[reports++] := in;
       }

       module
       report () {
           int(i);

           for (i=0 ; i<reports ; i++)
               .....  p_report[i];
       }
    \end{lstlisting}
    Each time module {\bf m} is invoked it establishes a reference
    in value array {\bf p\_report} to the module input queue. At the
    end of the module {\bf report} is invoked which iterates
    through the value array {\bf p\_report} and assigns the referenced
    queues to something. We might like to do something like the following
    \begin{lstlisting}[gobble=4, language=threepl]
       when (<p_report[0])
           p_out <<- p_report[0];
       else when (<p_report[1])
           p_out <<- p_report[1];
       else when (<p_report[2])
           p_out <<- p_report[2];
       .
       etc
       .
    \end{lstlisting}
    which will read from one of the queues in a priority order. We could
    generalise the above program by means of recursion - \index{recursion}
    \begin{lstlisting}[gobble=4, language=threepl]
           queue({result, q1, q2, q3}, "uint:8");
           value(v, "[3]uint:8");
           .
           .
           v[0] := q1;
           .
           .
           v[1] := q2;
           .
           .
           v[2] := q3;
           .
           .
           [pri(result <- v, 0, 3);]

       procedure
       pri (queue(o) <- value(i), int({index, size})) {
           when (<i[index])
               o <<- i[index];
           else
               if (index < (size - 1))
                   pri(o <- i, index+1, size);
       }
    \end{lstlisting}    
    All of the above code examples using nested {\em when}s will generate
    cascaded logic which can quickly lead to control logic paths which are
    too long for signals to propagate within a clock cycle.

    The use of immediate code to control target code generation can
    be illustrated by the following examples. If we wish to create
    a delay line we could write the following -
    \begin{lstlisting}[gobble=4, language=threepl]
           static({y, z}, "uint:8");

           .
           .
           z <- del(y, 5);
           .
           .

       function
       delay (value(in, "uint:8"), int(n)) {
           static(d, "[10]uint:8");
           [
               d[0] <- in;
               for (int(i,1) ; i<n ; i++)
                   d[i] <- d[i-1];
           ]
           return(d[n-1]);
       }
    \end{lstlisting}
    We have had to include the code for the delay line itself in a module,
    the function just returning a value which is the last static variable in
    the delay line. Thus the function itself does not contain any in-line
    target code (functions cannot), the module being an independent executing
    entity.
    
    The {\em for} loop is immediate and executes {\em n} times. Each time it
    iterates it generates one target assignment statement. Since we are at
    the top level of a module the target assignment statement before the loop
    and each of the target assignment statements generated by the loop will
    each be wrapped in their own infinite loops to run completely
    independently in parallel. Had we used a loop like this below the top
    level of a module it would have had to be enclosed in a {\em seq}, {\em
    par} or {\em sync}. The function above, as called, is equivalent to
    having written -
    \begin{lstlisting}[gobble=4, language=threepl]
       function
       delay (value(in, "uint:8"), int(n)) {
           static(d, "[10]uint:8");
           [
               d[0] <- in;
               d[1] <- d[0];
               d[2] <- d[1];
               d[3] <- d[2];
               d[4] <- d[3];
           ]
           return(d[4]);
       }
    \end{lstlisting}
    
    There are two annoying limitations of the above code: the data type is
    fixed and the static array in the function has a fixed dimension. The
    following version removes both these limitations -
    \begin{lstlisting}[gobble=4, language=threepl]
           static({y, z}, "uint:8");

           .
           .
           z <- del(y, 5);
           .
           .

       function
       delay (value(in), int(n)) {
           static(d, "["+n+"]"+type(in));
           [
               d[0] <- in;
               for (int(i,1) ; i<n ; i++)
                   d[i] <- d[i-1];
           ]
           return(d[n-1]);
       }
    \end{lstlisting}
    Parameter {\em in} is now created without a data type, which it will
    take from the matching argument. It now uses that data type by getting the
    string version by means of {\em type()} and prepending the required
    dimension in brackets from parameter {\em n}. The data type of the
    function return will be the same as {\em d}, which is the same as the
    input {\em in}.

    If an immediate loop is to contain more than one target statement
    then these must be bracketed, e.g.
    \begin{lstlisting}[gobble=4, language=threepl]
       for (i=0 ; i<3 ; i++) {
           a[i] <- b[i] + i;
           c[i] <- d[i] - i;
       }
    \end{lstlisting}
    is equivalent to writing
    \begin{lstlisting}[gobble=4, language=threepl]
       a[0] <- b[0] + 0;
       c[0] <- d[0] - 0;
       a[1] <- b[1] + 1;
       c[1] <- d[1] - 1;
       a[2] <- b[2] + 2;
       c[2] <- d[2] - 2;
    \end{lstlisting}    

    If we were to use such a loop inside a procedure we would have to enclose
    the resulting target statements in a {\em seq}, {\em par} or {\em sync}
    because the body of a procedure can only contain one top level target
    statement -
    \begin{lstlisting}[gobble=4, language=threepl]
       procedure
       p () {
           par {
               for (i=0 ; i<3 ; i++) {
                   a[i] <- b[i] + i;
                   c[i] <- d[i] - i;
               }
           }
       }
    \end{lstlisting}
    in this case a {\em par} block.

    An immediate loop that contains target code in its
    body simply replicates the target code and it is the context
    in which the immediate loop occurs that decides whether the
    resultant replicated statements are to be executed
    sequentially, in parallel or in individual infinite loops.
    For example array variables are not allowed in expressions, so if
    we wish to do array arithmetic we must code it explicitly. For example
    to add two 2D arrays we could write -
    \begin{lstlisting}[gobble=4, language=threepl]
       static({a, b, c}, "[3][4]uint:8");
       
       par {
           for (int(i) ; i<3 ; i++)
               for (int(j) ; j<4 ; j++)
                   c[i][j] <- a[i][j] + b[i][j];
       }
    \end{lstlisting}
    If this code is at the top level in a module then we may not want the {\em
    par}. However, note that this will execute on every clock cycle so we may
    want some execution control using a {\em when}. If the variables were
    queues rather than statics then their execution would be controlled by queue
    availability.
